\providecommand{\BM}{\textbf{[B]}}
\providecommand{\KM}{\textbf{[K]}}
\providecommand{\SM}{\textbf{[S]}}
\providecommand{\XM}{\textbf{[X]}}

\chapter*{Der Yukawa-Mechanismus als Konsequenz von $\tilde T\cdot m=1$}
\addcontentsline{toc}{chapter}{Der Yukawa-Mechanismus als Konsequenz von $\tilde T\cdot m=1$}

{\small\noindent\textbf{Zusammenfassung.}\quad\ignorespaces
Im Standardmodell ist die Yukawa-Kopplung $y_i = m_i/v$ eine empirische
Annahme: man setzt sie so an, dass nach der elektroschwachen Symmetriebrechung
die beobachteten Massen entstehen. Dieses Dokument zeigt, dass dieselbe
Kopplung aus der FFGFT-Grundrelation $\tilde T\cdot m = 1$ folgt, wenn man
die elektroschwache Skala $v$ als ortsabhängig zulässt. Die
massenproportionale Kopplung aller Fermionen ans Higgs ist dann keine
Annahme, sondern die fluktuierende Form von $\tilde T\cdot m = 1$. Das
Higgs-Teilchen $h(x)$ ist das Schwingungsquant der Gitterskala --- kein
eigenständiges Skalarfeld mit unbekanntem Potential. Seine Masse ist
die Reststeifigkeit des Gitters: $M_W^2+M_Z^2+m_h^2 = v^2/2$ (Spurregel,
$0{,}45\,\%$). Zusammen mit $\sin^2\theta_W = 2/9$ und $m_h/M_Z = 11/8$
folgen $M_Z$, $M_W$ und $m_h$ allein aus $v$ auf $0{,}2$--$0{,}3\,\%$.
}\medskip

\section*{Zeichenerklärung}
\addcontentsline{toc}{section}{Zeichenerklärung}

\medskip
\begin{tabular}{@{}ll@{}}
\textbf{[B]} & algebraisch bewiesen \\
\textbf{[K]} & numerisch bestätigt \\
\textbf{[S]} & strukturell plausibel, Herleitung offen \\
\textbf{[X]} & geprüft und negativ \\
\end{tabular}

\medskip
\begin{tabular}{@{}lp{9cm}@{}}
\toprule
\textbf{Symbol} & \textbf{Bedeutung} \\
\midrule
$\tilde T_i$ & Eigenzeit-Skalierung des Modus $i$ \\
$m_i$ & Ruhemasse von Teilchen $i$ \\
$v$ & elektroschwache Vakuumskala ($246{,}22$\,GeV) \\
$h(x)$ & physikalisches Higgs-Feld (Fluktuation von $v$) \\
$y_i$ & Yukawa-Kopplung (SM-Definition: $m_i = y_i v$) \\
$r_i$, $p_i$ & Vorfaktor und Exponent der FFGFT-Massenformel \\
$\xi$ & FFGFT-Skalierungsparameter ($4/30000$) \\
\midrule
$\psi_i$, $\bar\psi_i$ & Fermionfeld des Teilchens $i$ und sein Dirac-Adjungiertes \\
$\mathcal L$ & Lagrange-Dichte \\
$H$ & komplexe Gitteramplitude (Higgs-Feld), $H = (v+h)/\sqrt2$ \\
$n_{\phi,i}$, $n_{\theta,i}$ & Wicklungszahlen auf dem $T^4$-Torus \\
$n_c$ & Farb-Multiplizität \\
$d_{{\rm akt},i}$ & Zahl der aktiven Torus-Dimensionen (Dok.~189) \\
\midrule
$M_W$, $M_Z$ & Massen des $W$- und $Z$-Bosons (Polmassen) \\
$m_h$, $m_t$ & Masse des Higgs-Teilchens, des Top-Quarks \\
$m_s$ & Masse des Strange-Quarks \\
$W^\pm$ & geladenes Eichboson; eine komplexe Mode, $W^\pm = (W^1\mp iW^2)/\sqrt2$ \\
$W^3$, $B$ & neutrale Eichfelder vor der Mischung ($SU(2)$, $U(1)$) \\
$\gamma$ & Photon \\
$\theta_W$ & Weinberg-Winkel, on-shell $\sin^2\theta_W = 1 - M_W^2/M_Z^2$ \\
$g$, $g'$ & Eichkopplungen von $SU(2)$ und $U(1)$ \\
$\lambda$ & Higgs-Selbstkopplung, $m_h^2 = 2\lambda v^2$ \\
$y_t$ & Yukawa-Kopplung des Top-Quarks \\
$G_F$ & Fermi-Konstante \\
$\tau_\mu$ & Lebensdauer des Myons (einzige Messgröße hinter $v$) \\
\midrule
$\phi$ & goldener Schnitt $(1+\sqrt5)/2$ \\
$Q$ & Koide-Quotient $\sum m_i/(\sum\sqrt{m_i})^2$ \\
$\alpha$, $\alpha_s$ & Feinstrukturkonstante, starke Kopplung \\
$\Lambda_{\rm QCD}$ & Einschlussskala der starken Wechselwirkung \\
$R_{\rm Proton}$ & Ladungsradius des Protons \\
$\theta_{12}$, $\theta_{13}$, $\theta_{23}$ & Neutrino-Mischungswinkel \\
$\delta_{CP}$ & CP-verletzende Phase \\
$K_{\rm frak}$ & fraktaler Korrekturfaktor $1-100\xi$ (Dok.~133) \\
$\mathrm{GF}(q)$, $\mathrm{GF}(q)^*$ & Galois-Körper mit $q$ Elementen, seine multiplikative Gruppe \\
$A_5$ & Ikosaedergruppe (60 Elemente) \\
\bottomrule
\end{tabular}

\section{Die Grundrelation und ihre SM-Schreibweise}

Die FFGFT setzt für jeden Fermionmodus $i$:
\begin{equation}
  \tilde T_i \cdot m_i = 1, \qquad
  m_i = r_i\,\xi^{p_i}\,v.
  \label{eq:T0}
\end{equation}
Die Yukawa-Kopplung des Standardmodells ist definiert durch
$m_i = y_i\,v$, also:
\begin{equation}
  \boxed{y_i = r_i\,\xi^{p_i}.}
  \label{eq:yuk_ffgft}
\end{equation}
Das ist kein Matching, sondern eine direkte Umschreibung \BM.
Dabei ist (Dok.~189):
\begin{align}
  r_i &= \frac{n_{\phi,i}^2}{n_{\theta,i}\,n_c},
  & p_i &= \frac{d_{{\rm akt},i}}{3},
\end{align}
wobei $n_{\phi,i}$, $n_{\theta,i}$ die Wicklungszahlen auf dem $T^4$-Torus
und $n_c$ die Farb-Multiplizität sind. Die Yukawa-Kopplung ist damit
vollständig durch die Torus-Topologie bestimmt --- sie hat keinen freien
Parameter \BM.

\section{Ableitung des Yukawa-Vertex aus $\tilde T\cdot m = 1$}

\subsection{Der Vakuumfall}

Im Vakuum ist $v = {\rm const}$. Gl.~\eqref{eq:T0} gibt
$\tilde T_i = 1/m_i$ --- jeder Modus hat eine feste Eigenzeit-Skalierung.
Der Lagrange-Massenterm lautet:
\begin{equation}
  \mathcal L \supset -m_i\,\bar\psi_i\psi_i.
\end{equation}

\subsection{Fluktuierende Gitterskala}

Lasse nun zu, dass $v$ ortsabhängig fluktuiert:
\begin{equation}
  v(x) = v + h(x), \qquad |h(x)| \ll v.
\end{equation}
Die Grundrelation $\tilde T_i\cdot m_i = 1$ muss weiterhin gelten ---
nicht als globale Norm, sondern als lokale Bindung zwischen Eigenzeit und
Masse. Das bedeutet:
\begin{equation}
  m_i(x) = \frac{1}{\tilde T_i} = r_i\,\xi^{p_i}\,v(x)
  = r_i\,\xi^{p_i}\,\bigl(v + h(x)\bigr)
  = m_i^{(0)}\!\left(1 + \frac{h(x)}{v}\right).
  \label{eq:m_local}
\end{equation}

\subsection{Der Kopplungsterm}

Einsetzen in den Massenterm ergibt \BM:
\begin{equation}
  \mathcal L \supset -m_i(x)\,\bar\psi_i\psi_i
  = -m_i^{(0)}\,\bar\psi_i\psi_i
    \;-\; \frac{m_i^{(0)}}{v}\,\bar\psi_i\psi_i\,h(x).
  \label{eq:yukawa_derived}
\end{equation}
Der zweite Term ist exakt der Yukawa-Vertex des Standardmodells mit
$y_i = m_i^{(0)}/v$. Er folgt hier nicht aus einer Annahme über die
Kopplungsstruktur, sondern aus der einzigen FFGFT-Relation $\tilde T\cdot m = 1$
unter räumlichen Schwankungen von $v$.

\section{Massenproportionale Kopplung als Konsequenz}

Gl.~\eqref{eq:yukawa_derived} gilt für \emph{alle} Fermionen $i$
\emph{simultan}, mit demselben $v$ und demselben $h(x)$. Das bedeutet:

\begin{quote}
\textbf{[B]} Die massenproportionale Kopplung $y_i = m_i/v$ aller Fermionen
ans Higgs-Boson ist eine zwangsläufige Konsequenz von $\tilde T\cdot m = 1$
unter räumlicher Fluktuation der Gitterskala. Sie ist keine freie Annahme
des Modells.
\end{quote}

Das ist der Inhalt des Yukawa-Mechanismus in FFGFT-Sprache: nicht
„Teilchen koppeln an ein Skalarfeld", sondern „Masse ist eine lokale
Eigenschaft des Gitters, und $h(x)$ ist die Schwingung dieser Eigenschaft".

\section{Interpretation des Higgs-Teilchens}
\label{sec:interp}

Das Feld $h(x)$ ist in dieser Herleitung die ortsabhängige
Abweichung der Gitterskala $v$ vom Vakuumwert. Das physikalische
Higgs-Teilchen --- die Resonanz bei $125{,}2$\,GeV --- ist dann das
Schwingungsquant dieser Abweichung.

Das unterscheidet sich konzeptionell vom Standardmodell:

\begin{center}
\small\begin{tabular}{@{}lll@{}}
\toprule
 & \textbf{Standardmodell} & \textbf{FFGFT} \\
\midrule
Ursprung der Masse & Yukawa-Kopplung (angesetzt) & $\tilde T\cdot m = 1$ (Grundrelation) \\
Higgs-Feld & eigenständiges Skalarfeld & Fluktuation von $v(x)$ \\
Kopplung $y_i = m_i/v$ & empirischer Parameter & Ableitung aus Gl.~\eqref{eq:m_local} \\
Anzahl freier Parameter & $n_{\rm Fermionen}$ & $0$ \\
Higgs-Potenzial $V(H)$ & Mexikanerhut (angesetzt) & Reststeifigkeit des Gitters (§\ref{sec:spur}) \\
Higgs-Masse $m_h$ & freier Parameter & Spurregel (§\ref{sec:spur}) \KM \\
\bottomrule
\end{tabular}
\end{center}

Die Ableitung bis hierher zeigt, dass $h(x)$ existieren \emph{muss} und
mit $m_i/v$ koppeln \emph{muss}. Die Rückstellkraft der Schwingung --- das
Potential $V(h)$ und damit die Higgs-Masse --- behandeln §\ref{sec:higgsmasse}
und §\ref{sec:spur}.

\section{Higgs-Masse aus der Galois-Kette \KM}
\label{sec:higgsmasse}

Potential, Higgs-Masse und Schwingungsfrequenz sind ein einziges Problem.
Dieser Abschnitt zeigt,
dass ein numerisch scharfer Kandidat existiert, der vollständig aus
Galois-Strukturen des Korpus gebaut ist.

\paragraph{Die 11/8-Kette.}
Die drei Größen $M_Z$, $m_h$, $m_t$ erfüllen:
\begin{equation}
  M_Z : m_h : m_t \;=\; 1 : \tfrac{11}{8} : \bigl(\tfrac{11}{8}\bigr)^2,
  \qquad
  m_h = M_Z\cdot\tfrac{11}{8}, \quad m_t = M_Z\cdot\bigl(\tfrac{11}{8}\bigr)^2.
  \label{eq:kette}
\end{equation}
Numerisch \KM:
\begin{center}
\small\begin{tabular}{@{}lrrr@{}}
\toprule
Größe & gemessen & FFGFT (Gl.~\eqref{eq:kette}) & Abstand \\
\midrule
$M_W$ & $80{,}369$\,GeV & $M_Z\sqrt{7/9} = 80{,}420$\,GeV & $0{,}06\,\%$ \\
$M_Z$ & $91{,}188$\,GeV & Eingang & --- \\
$m_h$ & $125{,}200$\,GeV & $M_Z\cdot11/8 = 125{,}383$\,GeV & $0{,}15\,\%$ \\
$m_t$ & $172{,}570$\,GeV & $M_Z\cdot(11/8)^2 = 172{,}402$\,GeV & $0{,}10\,\%$ \\
\bottomrule
\end{tabular}
\end{center}
Alle vier schweren Bosonen/Fermionen folgen aus zwei Galois-Zahlen plus $v$.

\paragraph{Galois-Herkunft der Zahlen.}
Beide Faktoren sind Körpergrößen des Korpus:
$8 = |\mathrm{GF}(9)^*| = 3^2 - 1$ (Dok.~336) und
$11 \in \text{Orbit}\{7,11,21\}$ unter $x\mapsto 3x \bmod 26$
in $\mathrm{GF}(27)^*$ (Dok.~340, dieselbe Struktur wie $\Delta m^2_{\rm atm}$).
Der Weinberg-Faktor $\sqrt{7/9}$ kommt aus dem $A_5$-Matrixelement $2/9$
(Dok.~372). Die Kette verbindet damit drei verschiedene Galois-Strukturen
zu einem konsistenten Zahlenbild \KM.

\paragraph{Einordnung.}
Gl.~\eqref{eq:kette} ist eine numerisch bestätigte Relation; die
Galois-Herkunft der Zahlen 11 und 8 ist eine Lesart \SM. Eine Herleitung
der Higgs-Masse, die ohne diese Zahlen auskommt, liefert die Spurregel in
§\ref{sec:spur}; zusammen legen beide das gesamte elektroschwache
Bosonspektrum allein aus $v$ fest.

\section{Die Spurregel: Higgs-Masse als Reststeifigkeit des Gitters}
\label{sec:spur}

\paragraph{Bild.}
Gekoppelte Resonanzkreise verstimmen sich gegenseitig: die Kopplung schiebt
die Eigenfrequenzen auseinander. Was dabei erhalten bleibt, ist die Summe
der Frequenzquadrate --- die Spur der Steifigkeitsmatrix. Die Higgs-Masse
ist in diesem Bild keine eigene Resonanz mit eigenem Potential, sondern der
Anteil der Gittersteifigkeit, den die Eichbosonen übrig lassen.

\paragraph{Die Regel.}
\begin{equation}
  M_W^2 + M_Z^2 + m_h^2 \;=\; |\langle H\rangle|^2 \;=\; \frac{v^2}{2},
  \qquad
  m_h^2 = \frac{v^2}{2} - M_W^2 - M_Z^2 .
  \label{eq:spur}
\end{equation}

\paragraph{Die vier Bausteine.}
\begin{enumerate}
\item \textbf{Erhaltung unter Kopplung \BM.} Die Summe der Eigenwerte einer
  symmetrischen Matrix ist ihre Spur, unabhängig von der Basis. Die
  elektroschwache Mischung ist eine reine Drehung um $\theta_W$: vor der
  Mischung tragen $W^3$ und $B$ die Steifigkeiten $g^2v^2/4$ und
  $g'^2v^2/4$, danach Photon ($0$) und $Z$ ($M_Z^2$) --- die Summe ist
  identisch, die Mischmatrix orthogonal.
\item \textbf{Keine Eigensteifigkeit der Kopplung \BM.} Eine mechanische
  Koppelfeder ist ein eigenes Bauteil und legt Steifigkeit auf die
  Diagonale. Das Vakuum ist leer; die Kopplung kann nur umverteilen. Das ist
  exakt die Eigenschaft einer Drehung (Punkt~1).
\item \textbf{Zählung einmal pro Mode.} $W^+$ und $W^-$ sind eine komplexe
  Mode mit zwei Umlaufrichtungen: im Korpus ist das Antiteilchen dasselbe
  Feld mit umgekehrtem Vorzeichen (Dok.~042, 053), und die Spur läuft über
  $m^2$, wo das Vorzeichen herausfällt. Konsistent damit zählt Koide drei
  Leptonmassen, nicht sechs.
\item \textbf{Budget $= v^2/2$ \BM.} Das Higgs-\emph{Feld} ist das
  Grundgitter selbst; die Gesamtsteifigkeit gehört ihm und ist das
  Betragsquadrat seiner Vakuumamplitude. Da die Amplitude komplex ist
  (Betrag und Phase), verlangt die kanonische Normierung der reellen
  Schwingung $h$ den Faktor $1/\sqrt2$: $H = (v+h)/\sqrt2$, also
  $|\langle H\rangle|^2 = v^2/2$. Es ist dieselbe $\sqrt2$ wie in
  $W^\pm = (W^1\mp iW^2)/\sqrt2$ --- das Argument, das $W$ einmal zählen
  lässt, legt auch den Faktor $1/2$ fest.
\end{enumerate}
Das Higgs-\emph{Teilchen} $h$ ist die Eigenschwingung des Gitters; seine
Steifigkeit $m_h^2$ ist der Rest in Gl.~\eqref{eq:spur}. Das
Mexikanerhut-Potential des Standardmodells wird damit nicht angesetzt:
seine Krümmung am Minimum ist diese Reststeifigkeit.

\paragraph{Zahlen \KM.}
Mit $M_W^2 = \tfrac79 M_Z^2$ (Dok.~372) folgt
$m_h = \sqrt{v^2/2 - \tfrac{16}{9}M_Z^2} = 124{,}6$\,GeV
(gemessen $125{,}2$\,GeV, $-0{,}47\,\%$). Die Bilanz selbst stimmt
auf $0{,}45\,\%$.

\paragraph{Bosonspektrum allein aus $v$ \KM.}
Kombiniert man Gl.~\eqref{eq:spur} mit $\sin^2\theta_W = 2/9$ und der
Kette $m_h = \tfrac{11}{8}M_Z$ (§\ref{sec:higgsmasse}), so bleibt
\[
  M_Z^2\Bigl(\tfrac79 + 1 + \tfrac{121}{64}\Bigr) = \frac{v^2}{2}
  \quad\Longrightarrow\quad
  \frac{M_Z^2}{v^2} = \frac{288}{2113},
\]
und alle drei Bosonmassen folgen aus $v$ ohne gemessene Masse als Eingang:
\begin{center}
\small\begin{tabular}{@{}lrrr@{}}
\toprule
 & FFGFT & gemessen & Abstand \\
\midrule
$M_Z$ & $90{,}90$\,GeV & $91{,}19$\,GeV & $-0{,}31\,\%$ \\
$M_W$ & $80{,}17$\,GeV & $80{,}37$\,GeV & $-0{,}25\,\%$ \\
$m_h$ & $124{,}99$\,GeV & $125{,}20$\,GeV & $-0{,}17\,\%$ \\
\bottomrule
\end{tabular}
\end{center}
Als Kopplungen: $g = 0{,}6512$ (on-shell $2M_W/v = 0{,}6528$),
$\lambda = 0{,}1288$ (aus $m_h$: $0{,}1293$); die Higgs-Selbstkopplung ist
durch $4g^2/7 + 2\lambda = 1/2$ festgelegt und kein freier Parameter mehr.

\paragraph{Genauigkeit.}
Alle drei Werte liegen gleichsinnig $0{,}2$--$0{,}3\,\%$ zu tief.
Gl.~\eqref{eq:spur} ist eine Baumgraphen-Relation; die elektroschwachen
Strahlungskorrekturen liegen bei einigen Prozent. Genauer als die hier
erreichte Übereinstimmung ist ohne Schleifenrechnung nicht zu erwarten.
Die Wahl von $v$ verschiebt das Ergebnis: mit der SM-Definition
$246{,}22$\,GeV stimmt die Bilanz auf $+0{,}45\,\%$, mit den
FFGFT-Werten aus Dok.~006/R104 ($248{,}9$ bzw.\ $245{,}6$\,GeV) auf
$-1{,}7$ bzw.\ $+1{,}0\,\%$. Welches $v$ das richtige ist, ist die in
R104/R112 registrierte Frage (konventionsfrei nur über $\tau_\mu$), keine
Lücke der Spurregel.

\paragraph{Fermionseite: die Spurregel gilt nur für Gittermoden.}
Die naheliegende Übertragung ,,das Top trägt dasselbe Budget``,
$m_t^2 = v^2/2$ ($y_t = 1$), ist unverträglich mit der Kette: zusammen mit
Gl.~\eqref{eq:spur} verlangt sie $m_t/M_Z = \sqrt{2113/576} = 1{,}915$,
gemessen ist $1{,}8925$ ($+1{,}2\,\%$), während $(11/8)^2 = 1{,}8906$ auf
$0{,}10\,\%$ trifft. Die Lesung $y_t = 1$ wird verworfen \XM: sie hat keine
Begründung außer der Analogie, und sie ist die schlechtere. Dazu kommt, dass
die Top-Masse selbst schemaabhängig ist --- Polmasse $172{,}6$\,GeV,
$\overline{\rm MS}$-Masse $162{,}5$\,GeV, Spanne $5{,}8\,\%$, entsprechend
$y_t = 0{,}99$ bzw.\ $0{,}93$. Quarks sind keine freien asymptotischen Moden
(Dok.~372 §4); die Spurregel gilt für die Eigenschwingungen des Gitters
($W$, $Z$, $h$), nicht für eingeschlossene Fermionen. Für die Top-Masse gilt
die Kette $m_t = (11/8)^2M_Z$ \KM.

\paragraph{Rationale Verhältnisse gegen $v$-abhängige Relationen.}
Die Genauigkeit der Relationen ordnet sich danach, ob sie $v$ enthalten:
\begin{center}
\small\begin{tabular}{@{}llr@{}}
\toprule
Relation & enthält & Abstand \\
\midrule
$M_W^2/M_Z^2 = 7/9$ & nur Polmassen & $0{,}06\,\%$ \\
$m_t/M_Z = (11/8)^2$ & nur Polmassen & $0{,}10\,\%$ \\
$m_h/M_Z = 11/8$ & nur Polmassen & $0{,}15\,\%$ \\
$M_W^2+M_Z^2+m_h^2 = v^2/2$ & $v$ & $0{,}45\,\%$ \\
$m_t^2 = v^2/2$ (verworfen) & $v$ & $1{,}2\,\%$ \\
\bottomrule
\end{tabular}
\end{center}
Verhältnisse gemessener Polmassen sind konventionsfrei; dort treffen die
rationalen Galois-Zahlen am besten --- das ist zu erwarten, denn es sind
exakte Zahlen, keine Parameter. $v$ dagegen ist keine Messung, sondern die
SM-Definition $(\sqrt2\,G_F)^{-1/2}$, deren irrationale Faktoren aus der
SM-Formulierung stammen und die eine Theoriekorrektur von $-0{,}44\,\%$ aus
dem Myonzerfall trägt (Dok.~351, R112). Jede Relation mit $v$ erbt diese
Konventionen. Der Rest der Spurregel ($0{,}45\,\%$) liegt in der Größe
dieser Korrektur; dass er aus der Definition von $v$ stammt und nicht aus
der Regel, ist eine prüfbare Vermutung \SM.

\section{Bestandsaufnahme des Teilchenzoos (Stand 21.~September 2026)}
\label{sec:bestand}

Die folgende Übersicht fasst zusammen, was der Korpus nach Dok.~372 und
diesem Dokument für jeden Sektor liefert. Maßgeblich ist jeweils die beste
Herleitung; ältere, überholte Formeln sind im Register (Dok.~190) vermerkt.

\begin{center}
\small
\begin{tabular}{@{}>{\raggedright\arraybackslash}p{2.2cm}>{\raggedright\arraybackslash}p{6.6cm}>{\raggedright\arraybackslash}p{2.3cm}l@{}}
\toprule
\textbf{Sektor} & \textbf{Größe} & \textbf{Quelle} & \textbf{Status} \\
\midrule
Geladene & Massenverhältnisse $e:\mu:\tau$ aus dem Matrixelement $2/9$ & 292, 293, 367 & \BM/\KM \\
Leptonen & $\phi$-Skelett der Gewichte, $4/3 = 2Q$ & 368, 371 & \BM \\
\midrule
Neutrinos & Massen aus den $\mathrm{GF}(27)^*$-Orbits & 340 & \KM \\
 & $\nu_3$ Dirac, $\nu_{1,2}$ Majorana & 343 & \BM \\
 & $\sin^2\theta_{13} = (25\xi)^{2/3}$ & 340 & \KM \\
 & $\sin^2\theta_{23}\in\{4/9,5/9\}$, $|\sin^2\theta_{23}-\tfrac12| = \tfrac1{18}$ & 372 & \KM \\
 & $\theta_{12}$ (nur Größenordnung), $\delta_{CP}$ & 340 & \SM \\
 & Absolute Skala $m_\nu$ (Ansatz) & 007 & \SM \\
\midrule
Eichbosonen & Photon und acht Gluonen masselos & 320, 321, 339 & \BM \\
 & $M_W^2/M_Z^2 = 7/9$ ($\sin^2\theta_W = 2/9$ on-shell) & 372 & \KM \\
 & $M_Z^2/v^2 = 288/2113$ ($M_Z$ aus $v$ allein) & 373 & \KM \\
\midrule
Higgs & Yukawa-Kopplung $y_i = m_i/v$ aus $\tilde T\cdot m = 1$ & 373 & \BM \\
 & $m_h = \tfrac{11}{8}M_Z$ & 373 & \KM \\
 & Spurregel $M_W^2+M_Z^2+m_h^2 = v^2/2$, Selbstkopplung festgelegt & 373 & \KM \\
\midrule
Top & $m_t = (\tfrac{11}{8})^2 M_Z$ & 373 & \KM \\
\midrule
Übrige & Exponenten $p_i = d_{\rm akt}/3$ & 189 & \KM \\
Quarks & Vorfaktoren $r_i$ im Galois-Raster; Herleitung offen & 006, 358 & \BM/\SM \\
 & Ausrollen wie bei Leptonen ($Q = 2/3$) & 372 & \XM \\
\midrule
Hadronen & Confinement als topologische Randbedingung & 319 & \SM \\
 & Gell-Mann--Okubo-Aufspaltung konsistent mit $m_s$ & 372 & \KM \\
 & $R_{\rm Proton}$ strukturell eingeordnet & 372 & \SM \\
 & $\Lambda_{\rm QCD}$ aus der $\xi$-Leiter & 372 & \XM \\
\midrule
Kopplungen & $\alpha$ aus Galois-Ordnungen & 338 & \KM \\
 & $\alpha_s$ über den Galois-Lauf & 160, R78 & \KM \\
 & Cabibbo-Winkel & 372 & \SM \\
\midrule
Skala & $v$ aus dem Gitter & 006, R104 & \KM \\
 & konventionsfreier Test nur über $\tau_\mu$ & 351, R112 & \SM \\
\bottomrule
\end{tabular}
\end{center}

\paragraph{Ordnungsprinzip.}
Genau sind die freien Moden --- die geladenen Leptonen --- und die
Schwingungen des Gitters selbst ($W$, $Z$, $h$). Unscharf ist alles
Eingeschlossene: Quarks und Hadronen, deren Massen vom Schema abhängen.
Und die exakten Galois-Verhältnisse zwischen Polmassen treffen besser als
jede Relation, die die SM-Größe $v$ enthält (§\ref{sec:spur}).

\paragraph{Offen.}
$\theta_{12}$ und $\delta_{CP}$; die Quark-Vorfaktoren $r_i$ außer Top;
$\Lambda_{\rm QCD}$; der Cabibbo-Winkel; die geometrische Herleitung des
hadronischen Faktors in Dok.~160 (R123); die konventionsfreie Bestimmung
von $v$ über die Myon-Lebensdauer (R112).

\section{Ergebnis}

\begin{enumerate}
\item Die Yukawa-Kopplung $y_i = r_i\,\xi^{p_i}$ ist eine direkte
  Umschreibung der Torus-Topologie --- keine freie Annahme \BM.
\item Der Yukawa-Vertex $-(m_i/v)\,\bar\psi_i h\psi_i$ folgt aus
  $\tilde T\cdot m = 1$ unter $v\to v+h(x)$ für alle Fermionen
  simultan \BM.
\item Die massenproportionale Kopplung (Beweis-Fingerabdruck des Higgs)
  ist eine Konsequenz, nicht ein Parameter \BM.
\item $h(x)$ ist das Schwingungsquant der Gitterskala. Das Higgs-Potenzial
  ist keine eigene Setzung: seine Krümmung am Minimum ist die
  Reststeifigkeit des Gitters (§\ref{sec:spur}).
\item \textbf{Higgs-Masse-Kandidat} \KM: $m_h = M_Z\cdot11/8 = 125{,}383$\,GeV
  auf $0{,}15\,\%$; $m_t = M_Z\cdot(11/8)^2$ auf $0{,}10\,\%$. Beide
  Faktoren --- $11\in\mathrm{GF}(27)^*$, $8=|\mathrm{GF}(9)^*|$ --- sind
  Galois-Größen des Korpus; ihre Herkunft ist Lesart \SM.
\item \textbf{Spurregel}: $M_W^2+M_Z^2+m_h^2 = v^2/2$ auf $0{,}45\,\%$ \KM.
  Erhaltung unter Kopplung, fehlende Eigensteifigkeit, Zählung und Faktor
  $1/2$ begründet \BM; $m_h = \sqrt{v^2/2-\tfrac{16}{9}M_Z^2}$.
\item \textbf{Bosonspektrum aus $v$} \KM: $M_Z^2/v^2 = 288/2113$;
  $M_Z$, $M_W$, $m_h$ auf $0{,}17$--$0{,}31\,\%$, Baumgraphen-Niveau.
\item Bestandsaufnahme des gesamten Teilchenzoos in §\ref{sec:bestand}.
\item Top: Spurregel gilt nur für Gittermoden; $y_t = 1$ verworfen \XM,
  $m_t = (11/8)^2M_Z$ \KM.
\item Genauigkeit ordnet sich nach Konventionsabhängigkeit: rationale
  Polmassen-Verhältnisse $0{,}06$--$0{,}15\,\%$, Relationen mit $v$
  $0{,}45\,\%$. Offen \SM: ob der Rest der Spurregel aus der Definition von
  $v$ stammt (R104/R112).
\end{enumerate}

\medskip
\noindent\textbf{Prüfskript:} \texttt{2/python/Dok373\_Skripte/pruef\_373\_yukawa.py}

\medskip
\noindent\textbf{Register (Dok.~190):} R122 aktualisiert (Higgs-Masse: Kandidaten vorhanden).

\begin{thebibliography}{9}
\bibitem{dok006} J.~Pascher, \emph{Dok.~006 --- T0-Teilchenmassen} (2025).
\bibitem{dok189} J.~Pascher, \emph{Dok.~189 --- Torus-Ableitung der Wicklungszahlen} (2026).
\bibitem{dok372} J.~Pascher, \emph{Dok.~372 --- Das Matrixelement-Prinzip} (2026).
\end{thebibliography}
